\documentclass[conference]{IEEEtran}
\usepackage{cite}
\usepackage{amsmath,amssymb,amsfonts}
\usepackage{algorithmic}
\usepackage{graphicx}
\usepackage{textcomp}
\usepackage{xcolor}
\usepackage{hyperref}

\begin{document}

\title{Phishing Detection System Using Email Header Analysis}

\author{
\IEEEauthorblockN{1\textsuperscript{st} Tharun Chunchu}
\IEEEauthorblockA{\textit{Computer Engineering} \\
\textit{San José State University}\\
San Jose, USA \\
tharun.chunchu@sjsu.edu}
\and
\IEEEauthorblockN{2\textsuperscript{nd} [Student 2 Name]}
\IEEEauthorblockA{\textit{Computer Engineering} \\
\textit{San José State University}\\
San Jose, USA \\
email@sjsu.edu}
\and
\IEEEauthorblockN{3\textsuperscript{rd} [Student 3 Name]}
\IEEEauthorblockA{\textit{Computer Engineering} \\
\textit{San José State University}\\
San Jose, USA \\
email@sjsu.edu}
}

\maketitle

\begin{abstract}
As cyber threats continue to evolve, phishing remains one of the most prominent attack vectors compromising personal and corporate security. This paper presents a robust, Hybrid Phishing Detection System that merges a geometric heuristic metadata engine with an active Large Language Model (Google Gemini). By developing a Python-based FastAPI backend and a modern React web interface, the system securely ingests \texttt{.eml} files. The heuristic engine scrutinizes critical metadata fields including \textit{From}, \textit{Reply-To}, \textit{Message-ID}, and \textit{Received} routing hops. Simultaneously, the Gemini-2.5-Flash model evaluates the raw email payload for psychological manipulation. The results highlight the effectiveness of integrating both mathematical routing analysis and Natural Language Processing (NLP) to definitively mitigate social engineering attacks.
\end{abstract}

\begin{IEEEkeywords}
Phishing Detection, Email Security, Spoofing, Heuristics, \texttt{.eml} Parsing, FastAPI, React
\end{IEEEkeywords}

\section{Introduction}
\textit{Lead Author: Tharun Chunchu}

Phishing attacks have grown increasingly sophisticated, often bypassing standard email filters through advanced domain spoofing and routing manipulation. At the core of these attacks is the exploitation of the Simple Mail Transfer Protocol (SMTP), which traditionally lacks strict verification of sender identities. Consequently, malicious actors can forge email headers to impersonate legitimate entities, deceiving users into exposing sensitive information or executing malware. 

\subsection{Overview of the Problem and Motivation}
Despite the widespread deployment of security measures such as SPF (Sender Policy Framework) and DKIM (DomainKeys Identified Mail), spear-phishing remains a pervasive threat. Many clients fail to surface underlying header anomalies to end-users in a discernible manner. The motivation behind this project is to bridge the gap between complex analytical email forensics and accessible security tools. By building a transparent detection system, we allow users and organizations to manually verify suspicious correspondence effortlessly.

\subsection{Statement of Objectives}
The primary objective of this project is to develop and implement a full-stack Software Security application that identifies phishing emails. Specific goals include:
\begin{itemize}
    \item Creating an efficient Python-based parsing engine capable of extracting and sanitizing raw \texttt{.eml} metadata.
    \item Designing a heuristic scoring mechanism to detect spoofed domains, missing transmission IDs, and suspicious routing paths.
    \item Integrating the Google Gemini Generative AI Model to dynamically detect social engineering and coercive language in the email body.
    \item Developing an interactive frontend dashboard for seamless file upload and real-time security assessment.
\end{itemize}

\subsection{Report Outline}
The remainder of this report is organized as follows: Section II discusses the system architecture and technological stack. Section III details the implementation of the parsing logic and the heuristic detection algorithms. Section IV covers the frontend development and user interface integration. Finally, Section V concludes the report with a summary of the project's impact and findings.

\section{System Architecture}
\textit{Lead Author: [Student 2 Name]}

The architectural design of the Phishing Detection System is divided into a robust backend processing layer and a dynamic client-side interface, adhering to modern secure software development principles.

\subsection{Backend Infrastructure}
The backend is constructed using FastAPI, a modern, fast web framework for building APIs with Python. FastAPI was chosen for its high performance, asynchronous request handling, and automatic generation of security-compliant OpenAPI documentation. The backend exposes the \texttt{/api/scan} endpoint, which accepts multipart form data representing the uploaded \texttt{.eml} files.

\subsection{Frontend Infrastructure}
The frontend is built utilizing Vite and React. Instead of relying on vulnerable or bloated third-party template frameworks, the UI is styled utilizing bespoke CSS featuring modern visual aesthetics (such as glassmorphism and subtle micro-animations). This isolates the presentation layer entirely from the backend analytical engine, minimizing potential cross-site scripting (XSS) attack vectors by cleanly compartmentalizing data presentation rules.

\section{Implementation of the Hybrid Detection Engine}
\textit{Lead Author: [Student 3 Name] / Tharun Chunchu}

The core contribution of this project resides in the dual-layered hybrid analysis engine. The system pairs a mathematical structural evaluation of SMTP headers with an intelligent generative linguistic assessment.

\subsection{Email Header Parsing}
To safely unpack the email metadata, the system utilizes Python's native \texttt{email.parser.BytesParser}. This guarantees that the stream is evaluated securely without executing potentially malicious payloads embedded in HTML bodies. The engine targets specific metadata tags:
\begin{enumerate}
    \item \textbf{From and Reply-To}: Used to establish sender identity.
    \item \textbf{Subject}: Analyzed for social engineering hallmarks (e.g., ``Urgent'', ``Suspended'', ``Verify'').
    \item \textbf{Message-ID}: A universally unique identifier assigned by the originating mail server.
    \item \textbf{Received}: A sequential list of hops tracing the email's origin through specific IP addresses and MTAs (Mail Transfer Agents).
\end{enumerate}

\subsection{Heuristic Scoring Algorithm}
Once the metrics are cleanly extracted, the system runs a sequential heuristic evaluation. Points are accumulated to determine an active risk score:
\begin{itemize}
    \item \textbf{Domain Mismatch (Score +3)}: The most critical spoofing indicator. The engine extracts the top-level and secondary domains from both the \textit{From} and \textit{Reply-To} fields via regex isolation. Diverse or conflicting domains signal an immediate spear-phishing attempt.
    \item \textbf{Transmission Anomalies (Score +2)}: If an email lacks a \textit{Message-ID} or possesses an irregularly formatted ID, this heavily implies direct script injection rather than routing through a legitimate corporate outbound server.
    \item \textbf{Routing Shortfalls (Score +1)}: If fewer than two \textit{Received} headers exist, the email was likely executed via local bulk-mail tools directly to the target, bypassing standard external infrastructure.
\end{itemize}

\subsection{Active LLM Exploit Assessment}
To combat sophisticated spear-phishing attacks that bypass technical routing triggers by utilizing legitimately acquired domains, the system employs the Google Gemini (\texttt{gemini-2.5-flash}) Large Language Model. The raw message body is passed through a secure prompt funnel acting as an Elite Cybersecurity Analyst. If the LLM zero-shot classifier detects artificial urgency, explicit credential scraping, or dangerous manipulation within the text, the module forcefully injects a \textbf{+5 point penalty} to the baseline Threat Score. This aggressively overrides the technical categorization, guaranteeing that emails with authentic structural headers but malicious social-engineering bodies are properly halted and flagged as ``Phishing'' on the React dashboard.

Depending on the accumulated score threshold, the hybrid logic categorizes the file rigidly into safe, suspicious, or explicit phishing bounds.

\section{Results and Output Validation}
\textit{Lead Author: Tharun Chunchu / Shared}

To ensure validity, the system was aggressively tested against curated datasets. Using automated \texttt{pytest} pipelines, we ensured that boundary edge cases—such as malformed headers, blank lines, or Unicode anomalies—did not crash the parsing engine.

Furthermore, we simulated attacks using mock datasets derived from standard Open Source resources resembling the Enron Email dataset. When fed a benign message, the system cleanly traced identical domains and reported a \textit{Safe} classification. However, when presented with a fabricated attack originating from a legitimate un-spoofed domain containing an urgent credential-harvesting scam, the heuristic engine alone would have cleared it as technically safe. Thanks to the Hybrid engine logic, the Gemini LLM instantaneously detected the manipulation, elevated the score to maximum, and halted the user with explicit textual warnings detailing the psychological anomalies. 

\section{Conclusions}
\textit{Lead Author: Tharun Chunchu}

In this project, we successfully modeled, designed, and implemented a Hybrid Software Security system to combat email phishing and SMTP spoofing. By synthesizing deterministic header evaluation with the natural language comprehension of Google Gemini, we built a tool that provides immediate, foolproof feedback to users. The integration of Python's fast parsing capabilities alongside a secure React frontend ensures scalability and resilience. The findings demonstrate that examining technical transmission metadata alone is insufficient; fusing it with active Generative AI provides a truly comprehensive detection mechanism that effectively mitigates user susceptibility to sophisticated social engineering. 

\begin{thebibliography}{00}
\bibitem{b1} W. Stallings, \textit{Cryptography and Network Security: Principles and Practice}, 8th ed. Pearson, 2020.
\bibitem{b2} W. A. Conklin, G. White, \textit{Principles of Computer Security: CompTIA Security+ and Beyond}, 6th ed. McGraw Hill, 2022.
\bibitem{b3} A. Almomani, B. B. Gupta, S. Atawneh, A. Meulenberg, and E. Almomani, ``A Survey of Phishing Email Filtering Techniques,'' \textit{IEEE Communications Surveys \& Tutorials}, vol. 15, no. 4, pp. 2070--2090, 2013.
\bibitem{b4} J. Ma, L. K. Saul, S. Savage, and G. M. Voelker, ``Beyond blacklists: learning to detect malicious web sites from suspicious URLs,'' \textit{Proceedings of the 15th ACM SIGKDD international conference on Knowledge discovery and data mining}, 2009.
\end{thebibliography}

\end{document}
